% Generated by tools/edn_latex.py from reports/edn.md.
% Do not edit: change reports/edn.md and run `make edn`.
\ednentry{
  id = {EDN-55},
  title = {\texorpdfstring{\texttt{evaluate}}{evaluate} resumes the run \texorpdfstring{\texttt{train}}{train} created instead of opening its own},
  date = {2026-09-30},
  milestone = {M3: Quality Assurance},
  topic = {Experiment Tracking, Evaluation Protocol},
  participants = {@lukas2510, decided by the repository owner},
  decision = {\texttt{train@\allowbreak{}<variant>} creates the MLflow run and writes its id into \texttt{models/\allowbreak{}<variant>/\allowbreak{}model.\allowbreak{}json}. \texttt{evaluate} reads that id and reopens the same run with \texttt{resume\_\allowbreak{}run}, appending the test metrics and the gate verdict. One experiment holds exactly one run per variant, and no stage opens a run of its own except the stage that first produces one.},
  rationale = {\emph{Alternatives considered:}
\begin{itemize}[leftmargin=*, topsep=2pt]
\item \textbf{Option A: one run per DVC stage, joined by a shared tag.} This is what the course demo does and what its notes warn about.
\newline Pros: the simplest possible wiring; each stage owns what it logs and nothing is shared; no stage depends on another's identifier.
\newline Cons: four variants become eight runs, and the thing anyone actually wants -- the hyperparameters next to the MdAPE they produced -- is in two different rows that a reader has to join by hand. MLflow's compare view compares runs, not pairs of runs, so the one view the rubric asks for (``help you and others compare, understand and extend your results'') cannot be built from it. It also doubles again as soon as \#38 adds energy figures.
\item \textbf{Option B: nested runs, a parent per pipeline state with four children.}
\newline Pros: MLflow supports it natively and it expresses the ladder's structure.
\newline Cons: the parent has no natural owner. The four \texttt{train@\allowbreak{}<variant>} stages are independent DVC stages that can run in any order or alone, so whichever one created the parent would be doing it on behalf of the others, and \texttt{dvc repro train@\allowbreak{}lgbm-\allowbreak{}basic} on its own would either create a parent with one child or reuse a stale one. MLflow's compare view already handles four siblings, so the structure buys nothing it does not cost.
\item \textbf{Option C (chosen): \texttt{train} creates the run, \texttt{evaluate} resumes it.}
\newline Pros: one run per variant carries its hyperparameters, its train and validation loss, its model artefact, \#38's energy figures and \#39's test metrics and gate verdict, so the comparable view is the experiment's own table with no joining. It is also the natural place for the artefact, because the run that holds the metrics holds the model they were measured on.
\newline Cons: \texttt{evaluate} now depends on an identifier produced by another stage, so a model trained with tracking off has no run to append to and its metrics are dropped. That is visible rather than silent -- \texttt{model.\allowbreak{}json} says \texttt{tracking\_\allowbreak{}mode: \textquotedbl{}disabled\textquotedbl{}} and \texttt{run\_\allowbreak{}id: null}, and \texttt{resume\_\allowbreak{}run} logs which variant it could not resume -- and \texttt{RECOMMENDITOS\_\allowbreak{}REQUIRE\_\allowbreak{}TRACKING=\allowbreak{}1} makes it fatal for the runs that matter (EDN-54). A resumed run's recorded duration is also the training run's, not the evaluation's, so run duration in the UI means ``how long the fit took''.
\end{itemize}
\par
\emph{Why:} The demo's own notes flag one-run-per-stage as a wart, and the rubric asks for results that can be compared and understood. Resumption is the only one of the three options where the comparison needs no post-processing, and the cost it adds -- a cross-stage identifier -- is small and is carried in an artefact both stages already depend on.
\par
Two details were raised by the \#39 side of the interface and are part of the decision rather than of the implementation. \texttt{resume\_\allowbreak{}run} must be a context manager that ends the run it opened, because \texttt{evaluate} is one stage covering all four variants in one process and a run left open would collect the next variant's metrics. And a metric a stage could not measure is dropped rather than logged, because SC-04 to SC-06 are \texttt{null} until their machinery exists and MLflow has no representation for ``not measured'' -- so the absence of the key is the record, which is consistent with the metrics artefact recording \texttt{null} and never a fabricated pass.
\par
Verified against the team's real DagsHub server rather than only against a local store: four runs for four variants, each with 9 to 12 parameters, the fit metrics, the four tags this stage sets and the three provenance tags the seam sets, and the whole model bundle uploaded under \texttt{model/\allowbreak{}}. Resuming one of them appended the test metrics to it, left its status \texttt{FINISHED}, correctly skipped a \texttt{null} criterion, and \textbf{did not create a fifth run}. Filtering the experiment by \texttt{tags.\allowbreak{}git\_\allowbreak{}commit} returned exactly those four, which is the comparable view, and the filter string is written into the model card so a grader can reproduce it.
\par
One property of MLflow that the decision does not fix, and that the model card states: there is no upsert by run name, so a second \texttt{dvc repro} of one variant creates a second run rather than replacing the first. The earlier run stays as history, which is what tracking is for, and the view is filtered by the commit of the state being looked at.},
  ai = {Alternative generation, Alternative assessment, Recommendation, Solution generation},
  response = {Accepted},
  assessment = {The resumption design came out of a negotiation between the agents working on issues \#37 and \#39 rather than from one of them, which is also where the context-manager requirement came from -- the \#39 side pointed out that one stage covering four variants would otherwise leak a run across variants. AI then verified the result against the real server instead of asserting it, including the run count before and after resuming, which is the assertion that distinguishes this option from option A.},
  aievidence = {Claude Code sessions on 2026-09-30 designing and implementing issues \#37 and \#39: the seam was agreed between the two over three rounds, and the live DagsHub verification of the run count, the parameters, the tags and the uploaded artefacts was run in a throwaway experiment that was deleted afterwards.},
  evidence = {\href{https://github.com/mlops-2627q1-mds-upc/MLOPS_Recommenditos/blob/main/recommenditos/tracking.py}{\texttt{recommenditos/\allowbreak{}tracking.\allowbreak{}py}}, \texttt{resume\_\allowbreak{}run}; \href{https://github.com/mlops-2627q1-mds-upc/MLOPS_Recommenditos/blob/main/recommenditos/modeling/train.py}{\texttt{recommenditos/\allowbreak{}modeling/\allowbreak{}train.\allowbreak{}py}} and \href{https://github.com/mlops-2627q1-mds-upc/MLOPS_Recommenditos/blob/main/recommenditos/modeling/evaluate.py}{\texttt{recommenditos/\allowbreak{}modeling/\allowbreak{}evaluate.\allowbreak{}py}}; \texttt{tests/\allowbreak{}test\_\allowbreak{}model.\allowbreak{}py::test\_\allowbreak{}one\_\allowbreak{}variant\_\allowbreak{}is\_\allowbreak{}one\_\allowbreak{}run\_\allowbreak{}that\_\allowbreak{}evaluate\_\allowbreak{}appends\_\allowbreak{}to}; \href{https://github.com/mlops-2627q1-mds-upc/MLOPS_Recommenditos/blob/main/docs/docs/model-card.md}{model card}, Training Procedure, Experiment tracking; \hyperref[EDN-54]{EDN-54}; \href{https://github.com/mlops-2627q1-mds-upc/MLOPS_Recommenditos/issues/37}{issue \#37}; \href{https://github.com/mlops-2627q1-mds-upc/MLOPS_Recommenditos/issues/39}{issue \#39}.},
}
